% ============================================================
% CHAPTER: Conclusion
% Synthesizes all six contributions, answers the five RQs
% stated in Chapter 1, identifies cross-cutting findings,
% discusses limitations and future work. Rebalanced so no
% single contribution (QUEST/Chapter 6) dominates the synthesis
% sections.
% ============================================================

\chapter{Conclusion}
\label{chap:conclusion}

\startcontents[chapters]
\printmyminitoc{
\textit{Chapter summary.} This thesis pursued five research questions on making remote sensing information accessible through natural language, across six contributions spanning single-modality RSVQA robustness and architecture, multi-modal and multi-task extensions, and a complementary retrieval paradigm. Four of the five questions receive an affirmative answer; the fourth, whether jointly training RSVQA with other tasks measurably helps, does not, at the level of aggregate metrics, a negative result that itself constitutes a contribution. This chapter synthesizes these findings, identifies patterns that only become visible across contributions rather than within any single one, and discusses this thesis's limitations and the directions it leaves open.
}
% ============================================================

\section{Summary of Contributions}
\label{sec:conclusion_summary}

This section revisits each of the thesis's five research questions in turn, stating what the corresponding contribution actually found, rather than only what it attempted.

\textbf{RQ1 (Chapter~\ref{chap:llm_augmentation}): can data-centric augmentation improve RSVQA models' robustness to linguistic variation, without sacrificing in-distribution accuracy?} Yes. Generating syntactically diverse reformulations of training questions with an instruction-tuned LLM (LLaMA-2-7B), rather than relying on back translation, improved robustness to unseen phrasings substantially: on questions reformulated by the LLM itself, the augmented model retained 81.77\% accuracy, close to its 82.98\% in-distribution accuracy, while the non-augmented model dropped to 52.16\% on the same reformulations. In-distribution accuracy was preserved (82.98\% versus 82.17\% for the non-augmented model), unlike back translation, which measurably degraded it (69.07\%), indicating that the syntactic diversity LLM-generated reformulations offer, rather than back translation's more limited variation, is what drives the generalization gain.

\textbf{RQ2 (Chapter~\ref{chap:seg_attention}): does conditioning attention on an explicit semantic segmentation prior improve RSVQA over attention computed from raw visual features alone?} Yes. Guiding attention with a frozen, multi-channel segmentation prior improved overall accuracy from 36.26\% (no attention) to 41.43\% (vanilla attention) to 45.44\% (segmentation-guided attention), a further gain of roughly 4 points from segmentation guidance specifically, on top of what attention alone already provides. This indicates that giving the model an explicit, separately trained signal for scene structure is more effective than requiring it to discover that structure from question-answering supervision alone.

\textbf{RQ3 (Chapter~\ref{chap:tammi}): how can RSVQA be extended to jointly exploit multiple sensing modalities acquired at different spatial resolutions?} By treating the resolution mismatch between modalities as a deliberate, user-configurable design choice rather than a limitation to be resolved architecturally: TAMMI extracts multispectral and SAR context at an independently configurable spatial extent around each VHR patch, and MM-RSVQA fuses the resulting three-modality representation with the question through trainable attention. Multimodal fusion improved average accuracy from 61.94\% (VHR-only) to 65.56\% (all three modalities), with the largest gains concentrated in questions requiring context beyond what a single VHR patch's extent can contain, such as urban classification (+24.6 points) and administrative department (+9.2 points). This benefit is not solely a matter of broader spatial context, however: SAR's contribution peaked specifically on land cover classification (74.89\%, the highest accuracy across every configuration tested, including the full three-modality model), consistent with SAR's physical sensitivity to surface roughness and dielectric properties rather than with the additional spatial extent it provides. Modalities can therefore be complementary along two distinct axes, the spatial context they cover and the physical property they are sensitive to, and both, not context alone, account for multimodal fusion's benefit here.

\textbf{RQ4 (Chapter~\ref{chap:multitask}): does training RSVQA jointly with captioning and self-supervised image reconstruction measurably improve either task?} No, not at the level of aggregate metrics. Single-task training matched or outperformed every multi-task configuration on both captioning and VQA. A closer look, however, qualified this result rather than simply confirming it: a modality-specific weakness in the multispectral reconstruction plausibly diluted the aggregate reconstruction signal, and a qualitative counter-example showed the full multi-task model succeeding on a specific question where the single-task baseline failed, meaning the aggregate result does not rule out category-specific synergy, even though it does not establish it either.

\textbf{RQ5 (Chapter~\ref{chap:ctcr}): can a composed query, combining a visual reference and a textual description, support retrieval of temporal change patterns in Earth Observation archives?} Yes. Composed Temporal Change Retrieval, introduced as a task in this thesis, is addressed effectively by decomposing the composed query into its reference and textual components before fusion (StageCTCR), which outperformed a direct end-to-end fusion of the two (UniCTCR). This advantage is not attributable to the staged decomposition alone: StageCTCR also leverages a pretrained vision-language model that UniCTCR's more directly fused architecture does not benefit from in the same way, so the two factors are confounded in this comparison rather than isolated. What can be said is that staged decomposition combined with pretrained vision-language representations outperformed direct end-to-end fusion; disentangling how much each factor individually contributes is left to future work.

\section{Cross-Cutting Findings}
\label{sec:conclusion_crosscutting}

Considered individually, each contribution answers the research question it was designed to address. Considered together, three patterns emerge that are not visible from any single chapter alone.

\textbf{RSVQA's remaining weaknesses are distributed across several independent levers, not concentrated in one.} Chapters~\ref{chap:llm_augmentation} and~\ref{chap:seg_attention} improve RSVQA from two different angles, data-centric augmentation and architectural attention design, without either depending on the other; Chapter~\ref{chap:tammi} then adds a third, largely orthogonal lever, showing that modality choice itself, independent of both linguistic diversity and attention design, measurably changes what a model can answer, and does so through two distinct mechanisms, spatial context and sensor-specific physical sensitivity, rather than one. Taken together, these three contributions suggest that RSVQA's remaining gaps are not a single bottleneck a lone intervention could close, but several largely independent axes, each requiring its own targeted solution.

\textbf{Reconciling heterogeneous visual signals is a recurring design problem, addressed differently each time it recurs.} Chapter~\ref{chap:seg_attention} reconciles a question with a segmentation prior at the same spatial resolution; Chapter~\ref{chap:tammi} reconciles three modalities at three different resolutions by making that mismatch a configurable dataset property rather than an architectural constraint; Chapter~\ref{chap:multitask} reconciles those same modalities with two additional training objectives within one architecture; and Chapter~\ref{chap:ctcr} reconciles a visual reference with a textual description of change within a composed query. No single solution recurs across these four instances, an attention mechanism, a dataset design choice, a shared multi-objective architecture, and a staged retrieval pipeline, respectively, suggesting that reconciling heterogeneous visual and textual signals in remote sensing is not a problem this thesis's contributions solve once and reuse, but one that resurfaces in a different form at each new combination of modality, task, and resolution.

\textbf{Assumed multi-task synergy does not straightforwardly hold.} Chapter~\ref{chap:relatedworks} noted that RS-LLaVA's own reported results show single-task fine-tuning outperforming joint training on both captioning and VQA metrics, an isolated observation drawn from prior work rather than a result this thesis produced itself. Chapter~\ref{chap:multitask} then measured the same question directly, on a different dataset, a different architecture, and a different task combination, and found the same pattern. Two independent measurements of the same negative result is a stronger basis for caution than either alone, and should be treated as reason to test rather than assume similar synergy claims elsewhere, including Chapter~\ref{chap:ctcr}'s own finding, where the respective contributions of staged decomposition and pretrained representations were deliberately not assumed but left as an open, explicitly confounded question rather than a settled one.

\section{Limitations}
\label{sec:conclusion_limitations}

Beyond the limitations already discussed within each individual chapter, three limitations apply across this thesis as a whole.

\textbf{Geographic and sensor scope.} The datasets introduced in this thesis, TAMMI and its extensions, are constructed over Metropolitan France specifically, combining VHR, Sentinel-2, and Sentinel-1 imagery. Whether the findings in Chapters~\ref{chap:tammi} and~\ref{chap:multitask} generalize to other geographic regions, land cover distributions, or sensor combinations was not tested, and is not something this thesis's results can be assumed to predict.

\textbf{Scale.} Each contribution in this thesis evaluates one architecture, at a scale consistent with academic compute constraints, rather than a systematic sweep across model or data scale. Whether the specific quantitative findings reported throughout this thesis, including RQ4's negative result on multi-task synergy, hold as model and data scale increase remains untested.

\textbf{Confounded comparisons left undisentangled.} Two contributions identify an effect without fully isolating its cause. In Chapter~\ref{chap:multitask}, the multispectral reconstruction weakness was localized to a specific, patch-internal artifact but not conclusively confirmed, and the per-question-type breakdown needed to determine whether multi-task synergy might hold for specific question categories despite the aggregate negative result was left incomplete. In Chapter~\ref{chap:ctcr}, StageCTCR's advantage over UniCTCR was shown to depend on both staged decomposition and pretrained vision-language representations jointly, without disentangling how much each factor individually contributes. Both cases identify a real effect while leaving its precise mechanism for future work to establish.

\section{Future Work}
\label{sec:conclusion_future}

The limitations above point directly to concrete extensions of this work, alongside broader directions this thesis's findings motivate but does not itself pursue.

Within individual chapters, two disentangling studies would directly extend this thesis's confounded findings. Completing the per-question-type breakdown for Chapter~\ref{chap:multitask}'s four training configurations would test whether the aggregate absence of multi-task synergy masks category-specific effects of the kind its qualitative example suggests, or holds uniformly; addressing the multispectral reconstruction weakness identified in the same chapter, most plausibly by replacing the affected decoder's index-based unpooling with a learned upsampling stage, would clarify whether it was a genuine contributor to reconstruction's limited measured benefit. Separately, evaluating an end-to-end CTCR-RS architecture that also benefits from pretrained vision-language representations would isolate how much of StageCTCR's advantage over UniCTCR, reported in Chapter~\ref{chap:ctcr}, is attributable to staged decomposition specifically, rather than to the two factors' current confound.

More broadly, this thesis's two directions, RSVQA and CTCR-RS's composed retrieval, were developed as complementary but separate paradigms. Whether the multi-modal, multi-resolution representations developed for RSVQA in Chapters~\ref{chap:tammi} and~\ref{chap:multitask} could support composed retrieval directly, or whether retrieval requires representations shaped differently from the ones question answering benefits from, is an open question this thesis's separation of the two paradigms does not itself answer. Finally, extending the datasets introduced in this thesis beyond Metropolitan France for TAMMI, and to further change captioning datasets for CTCR-RS, would test how far the specific findings reported here actually generalize, rather than assuming they do.

\section{Closing Remarks}
\label{sec:conclusion_closing}

This thesis set out to make remote sensing information more accessible through natural language, along two directions: answering questions about a known location, and retrieving locations that match a description. Across its five research questions, targeted interventions, data-centric augmentation, segmentation-guided attention, multi-modal fusion, and composed retrieval, each measurably improved on the specific gap they were designed to close, though rarely through a single mechanism: modality choice helped through two distinct physical properties rather than one, and composed retrieval's own gain was shown to rest on two confounded factors rather than a single clean cause. The one question that did not confirm what it set out to test, whether jointly training RSVQA with other tasks helps, is not a shortfall among these results but one of the more informative ones: an assumption common enough in the field to motivate real datasets and real model designs was tested directly, on new data and a new architecture, and did not hold, echoing an independent finding already present in the literature. Taken together, these six contributions leave remote sensing information more accessible through language than this thesis found it, and leave a clearer sense than before of which assumptions in that pursuit still need testing, rather than taking for granted.